% =============================================================================
% Hidden Currents: A Peak-Stress Water Index for Auditing AI Data Center
% Sustainability and Informing Mandatory Disclosure Policy
%
% v2 -- Revised in response to professorial peer review.
% Addresses reviewer points #1, #2, #3, #4, #5, #7, #8, #10, #13.
% =============================================================================

\documentclass[conference,10pt]{IEEEtran}
\usepackage[T1]{fontenc}
\usepackage[utf8]{inputenc}
\usepackage{amsmath, amssymb, amsfonts, amsthm}
\usepackage{graphicx}
\usepackage{booktabs}
\usepackage{cite}
\usepackage{url}
\usepackage{hyperref}
\usepackage{xcolor}
\usepackage{tcolorbox}
\usepackage[capitalise]{cleveref}

\hypersetup{colorlinks=true, linkcolor=blue!60!black,
            citecolor=blue!60!black, urlcolor=blue!60!black,
            pdftitle={Hidden Currents PSWI Paper}}

\newtheorem{proposition}{Proposition}

\begin{document}

\title{Hidden Currents: A Peak-Stress Water Index for Auditing AI Data Center Sustainability and Informing Mandatory Disclosure Policy}

\author{
\IEEEauthorblockN{Hrihaan Bhutani}
\IEEEauthorblockA{Emerald High School \\
Email: hribhu19@gmail.com \\
Code and data: \url{https://github.com/hrihaan19/pswi-data-center-audit}}
}

\maketitle

% ============================================================================
\begin{abstract}
% ============================================================================
The water consumption of data centers has emerged as a critical
sustainability concern, yet the metrics used to assess and report this
consumption obscure two of its most important dimensions: temporal
concentration (peak demand) and spatial heterogeneity (local water stress).
The industry-standard Water Usage Effectiveness (WUE) metric is reported as
an annualized fleet-wide average, treating every liter of water as
ecologically equivalent. Recent research has separately quantified the
peak-demand problem (Han et al. 2026) and the water-stress problem (Wu et
al. 2025), but no existing metric integrates both. This paper introduces
the \emph{Peak-Stress Water Index} (PSWI), a composite metric that
multiplies a facility's peak Water Usage Effectiveness (pWUE) by its
location-specific Baseline Water Stress (BWS) value from the WRI Aqueduct
4.0 dataset. We apply PSWI to a publicly auditable dataset of 53 hyperscale
and major colocation data centers operated by eight providers across North
America, Europe, and Asia. Our findings reveal a 4{,}100$\times$ disparity
between the best and worst facilities; a same-operator case study (Google
Lenoir, NC vs.\ Google Henderson, NV) demonstrates a 17$\times$ PSWI ratio
between facilities with similar self-reported WUE. A variance decomposition
shows that location-specific water stress accounts for roughly twice as
much cross-facility variation as efficiency does, and that covariance
terms among the three components account for nearly half of total
variation---an empirical justification for the multiplicative functional
form. Sensitivity analyses establish that facility rankings are robust to
$\pm$50\% per-facility perturbation of the peaking factor (mean Spearman
$\rho = 0.982$). A Monte Carlo counterfactual that propagates documented
uncertainty in WUE ($\pm$15\%), BWS ($\pm$25\%), and $\beta$ ($\pm$30\%)
shows that reallocating 10\% of fleet workload from worst-PSWI to
best-PSWI facilities yields a fleet-wide impact reduction of 36.7\% (95\%
CI $[26.1\%, 47.0\%]$). We conclude with a disclosure framework, modeled
on SEC Regulation S-K, that would mandate quarterly per-facility PSWI
reporting. All data and code are released under permissive licenses.
\end{abstract}

\begin{IEEEkeywords}
data centers, sustainable computing, water sustainability,
environmental policy, water stress, AI infrastructure, disclosure,
sensitivity analysis
\end{IEEEkeywords}

% ============================================================================
\section{Introduction}
\label{sec:introduction}
% ============================================================================

The rapid expansion of artificial intelligence has been accompanied by an
equally rapid expansion of the data center infrastructure that powers it.
Recent projections place U.S.\ data center electricity consumption at
6.7\% to 12.0\% of national totals by 2028~\cite{lbnl2024}, and the
International Energy Agency expects global data center electricity demand
to grow at double-digit annual rates through 2030~\cite{iea2025}. Less
widely discussed, but increasingly urgent, is the parallel growth in data
center \emph{water} demand. Between 2024 and 2030, U.S.\ hyperscale data
centers alone are projected to require an additional 697 to 1{,}451 million
gallons per day of water capacity---an amount comparable to the daily
potable supply of New York City~\cite{han2026}. The total capital cost of
expanding public water systems to meet this demand has been estimated at
\$10 to \$58 billion~\cite{han2026}. These numbers, however, hide a deeper
problem.

Industry's primary water sustainability metric, Water Usage Effectiveness
(WUE)~\cite{li2023thirsty}, treats every liter of consumed water as
equivalent. WUE is reported as an \emph{annualized} ratio of facility-level
water consumption to IT energy, typically averaged over an entire
year~\cite{google2024envreport, microsoft2024sustain, meta2025sustain}.
This treatment obscures two distinct dimensions of impact:

\begin{enumerate}
    \item \textbf{Temporal concentration of demand (peak load).} Han et
        al.~\cite{han2026} demonstrate that data center water use is
        highly concentrated on the hottest days of the year, with daily
        peaking factors routinely reaching 6 to 10. A facility may report
        a low annual WUE while simultaneously placing extreme stress on
        local water systems during summer peaks.

    \item \textbf{Spatial heterogeneity of stress (location).} The
        ecological cost of one liter of water consumed in Phoenix,
        Arizona---in the 22nd year of a megadrought affecting the lower
        Colorado River basin---is not equivalent to the cost of one liter
        consumed in Hamina, Finland, where sea-water cooling exploits
        effectively unlimited supply. Wu et al.~\cite{wu2025scarf}
        formalize this insight in their SCARF framework, weighting
        computing's water consumption by a location-specific stress factor
        derived from the World Resources Institute's Aqueduct
        dataset~\cite{kuzma2023aqueduct}.
\end{enumerate}

These two dimensions have been studied independently. Han et al.\ quantify
peak demand but treat all water as equivalent. Wu et al.\ weight by stress
but use annualized water totals. \emph{No existing metric combines both.}

\subsection{Contributions}

This paper makes five contributions toward closing this gap.

\paragraph{(1) A composite metric with empirical justification.}
We introduce the \emph{Peak-Stress Water Index} (PSWI), defined as the
product of peaking-adjusted Water Usage Effectiveness (pWUE) and
location-specific Baseline Water Stress (BWS). Unlike prior parallel
work, we ground the multiplicative functional form in an empirical
variance decomposition: covariance terms among the three components
account for 46.7\% of cross-facility variation in $\log\textrm{PSWI}$,
indicating that the components compound rather than add
(Section~\ref{sec:metric}).

\paragraph{(2) An open, auditable dataset.}
We assemble the first publicly available cross-company dataset of
hyperscale and major colocation data center water sustainability data,
covering 53 facilities operated by eight major providers across North
America, Europe, and Asia. Inclusion criteria are explicit and the sample
covers the majority of disclosed-WUE U.S.\ hyperscale capacity reported
in LBNL's 2024 census~\cite{lbnl2024} (Section~\ref{sec:methodology}).

\paragraph{(3) Robustness-tested empirical findings.}
We document a 4{,}100$\times$ disparity between best and worst
facilities. A same-operator case study isolates location effects from
methodology effects: Google's Lenoir, NC and Henderson, NV facilities
report similar WUE but differ in PSWI by a factor of 17. Bootstrap
resampling and Monte Carlo perturbation establish that these findings
are statistically robust. Sensitivity analyses show that facility
rankings remain stable under $\pm$50\% per-facility perturbation of the
peaking factor (mean Spearman $\rho = 0.982$); under $\pm$25\%
perturbation of BWS; and after inclusion of estimated Scope~2 water
($\rho = 0.946$, top-10 overlap 7/10) (Section~\ref{sec:results}).

\paragraph{(4) An uncertainty-propagated counterfactual.}
A Monte Carlo simulation that propagates documented relative
uncertainty in all three PSWI inputs shows that reallocating 10\% of
fleet workload from worst-PSWI to best-PSWI facilities yields a fleet
impact reduction of 36.7\% (95\% CI $[26.1\%, 47.0\%]$)---a substantial
margin even at the lower confidence bound (Section~\ref{sec:policy}).

\paragraph{(5) A policy framework.}
We propose a regulatory framework, modeled on SEC Regulation S-K, that
would mandate quarterly per-facility PSWI disclosure. We explicitly
engage with how the SEC's actual climate disclosure rules
(finalized March 2024 and partially stayed) handled comparable data
quality and auditability problems (Section~\ref{sec:policy}).

% ============================================================================
\section{Background}
\label{sec:background}
% ============================================================================

\subsection{Data Center Water Use}

Data centers consume water both directly and
indirectly~\cite{li2023thirsty, han2026}. Direct (Scope~1) water use is
dominated by on-site cooling. The high power density of modern GPU-based
AI servers, often exceeding 30~kW per rack~\cite{lbnl2024}, produces
heat loads that must be rejected to the outside environment. Two cooling
architectures predominate: \emph{evaporative cooling}, in which water is
intentionally evaporated to remove heat at high thermodynamic
efficiency, and \emph{dry cooling}, which uses mechanical refrigeration
or air-side economizers and consumes substantially more electricity.
Industry disclosures indicate water-cooled facilities can use 25--35\%
less electricity than air-cooled equivalents during peak summer
load~\cite{han2026, google2025envreport}. Indirect (Scope~2) water use
arises from upstream electricity generation, primarily thermoelectric
plants. While Scope~2 water typically exceeds Scope~1 in
volume~\cite{li2023thirsty}, it is rarely sourced from potable
infrastructure; we test sensitivity to its inclusion in
Section~\ref{sec:results}.

\subsection{The WUE Metric and Its Limitations}

Water Usage Effectiveness (WUE) is defined as:
\begin{equation}
\textrm{WUE} = \frac{\textrm{On-site annual water consumption (L)}}
                    {\textrm{Annual IT energy (kWh)}}
\label{eq:wue}
\end{equation}
WUE has three structural limitations that compromise its policy relevance.
\textbf{Annualization erases peak signal:} public water systems are sized
for peak daily demand, not annual averages. \textbf{Location-blindness
ignores ecological context:} a liter in a megadrought-affected basin is
treated identically to a liter on the Baltic. \textbf{Aggregation obscures
facility-level harm:} most operators report fleet-wide
averages~\cite{microsoft2024sustain, meta2025sustain, amazon2024aws},
diluting high-impact sites against low-impact ones.

% ============================================================================
\section{Related Work}
\label{sec:related}
% ============================================================================

We synthesize three previously parallel research streams.

\paragraph{Sustainable AI and computing.}
Strubell et al.~\cite{strubell2019energy} documented the energy footprint
of large language model training; Patterson et al.~\cite{patterson2022carbon}
extended this to carbon. Wu et al.~\cite{wu2022sustainable} provided an
industry perspective from Meta. Schwartz et al.~\cite{schwartz2020green}
introduced the Green AI framing. Luccioni et al.~\cite{luccioni2023bloom}
performed a life-cycle carbon analysis of BLOOM. Most of this literature
treats water as either secondary to carbon or omits it.

\paragraph{Geographical load balancing for sustainability.}
Islam et al.~\cite{islam2018water} pioneered water-aware load balancing.
Acun et al.~\cite{acun2023carbon} introduced Carbon Explorer.
Radovanovi\'{c} et al.~\cite{radovanovic2023carbon} described Google's
production carbon-aware system. Hanafy et al.~\cite{hanafy2024going}
quantified the cost premium of carbon reduction. Li et
al.~\cite{li2024equitable} introduced equity-aware GLB. Liu et
al.~\cite{liu2024gsr} proposed Geographical Server Relocation. Bian et
al.~\cite{bian2024cafe} applied carbon-aware federated learning. Yang
et al.~\cite{yang2023acrl} provided theoretical foundations. Across
this literature, water is treated through unweighted volumetric metrics.

\paragraph{Water-aware computing.}
Wu et al.'s SCARF framework~\cite{wu2025scarf} introduces the Adjusted
Water Impact (AWI) metric, multiplying volumetric consumption by a
location-specific stress factor from WRI
Aqueduct~\cite{kuzma2023aqueduct}. Jiang et al.'s
WaterWise~\cite{jiang2025waterwise} co-optimizes carbon and water in
scheduling; their ThirstyFLOPS~\cite{jiang2025thirsty} provides
region-specific HPC modeling. Han et al.~\cite{han2026} provide the
foundational quantitative analysis of data centers' impact on U.S.
public water systems and propose pWUE and Water Capacity Impact (WCI).
The de Vries-Gao study~\cite{devries2026} provides global AI water
estimates. Karimi et al.~\cite{karimi2022water} examine hot-arid
trade-offs. Siddik et al.~\cite{siddik2021environmental} document U.S.
data center environmental footprints.

\paragraph{Our contribution.}
The two structural deficiencies of WUE---peak blindness and
location-blindness---have each been addressed in isolation: Han et al.\
address peaks but treat all water as equivalent; Wu et al.\ address
location stress but use annual totals. \emph{No prior work integrates
the two.} We bridge this gap with PSWI, validate the multiplicative form
empirically through variance decomposition, and provide the first
publicly auditable cross-company dataset.

% ============================================================================
\section{The PSWI Metric}
\label{sec:metric}
% ============================================================================

\subsection{Definition}

We define the Peak-Stress Water Index for a data center facility $i$ as:
\begin{equation}
\textrm{PSWI}_i \;=\; \textrm{pWUE}_i \;\cdot\; \textrm{BWS}_i
\;=\; \textrm{WUE}_i \cdot \beta_i \cdot \textrm{BWS}_i
\label{eq:pswi}
\end{equation}
where $\textrm{WUE}_i$ is the annualized WUE in L/kWh; $\beta_i \geq 1$
is the daily peaking factor (the ratio of peak-day to average-day water
consumption, following Han et al.~\cite{han2026}); and
$\textrm{BWS}_i$ is the dimensionless ratio of total water demand to
renewable supply at the facility's catchment, drawn from WRI Aqueduct
4.0~\cite{kuzma2023aqueduct}. PSWI has units of (L/kWh) and admits the
interpretation: \emph{the volume of water a facility consumes per kWh
of IT energy on a peak day, weighted by the relative scarcity of that
water at its location.}

\subsection{Justification of the Multiplicative Form}

A reasonable question is whether the components should combine
multiplicatively (as in our definition) or additively. We argue both
theoretically and empirically that multiplication is the correct
operator.

\paragraph{Theoretical argument.}
Stressed water systems exhibit threshold behavior: when supply is
abundant, an additional liter of consumption produces marginal harm; as
the demand-to-supply ratio approaches and exceeds unity, an additional
liter produces nonlinearly greater harm because the system has no buffer
capacity. Similarly, peak consumption matters more in stressed systems
than in unstressed ones: a Phoenix facility's peak draws compete with
peak residential demand in a basin that already cannot meet baseline
demand. The two effects compound: peak consumption in a stressed region
produces nonlinearly greater harm than the sum of either effect in
isolation. A multiplicative form captures this compounding; an additive
form does not.

\paragraph{Empirical validation.}
A formal variance decomposition of $\log\textrm{PSWI}$ across our 53
facilities supports the multiplicative form. Since
$\log\textrm{PSWI} = \log\textrm{WUE} + \log\beta + \log\textrm{BWS}$,
the cross-facility variance decomposes as:
\begin{equation}
\textrm{Var}(\log\textrm{PSWI}) =
\sum_k \textrm{Var}(\log X_k) + 2 \sum_{j<k} \textrm{Cov}(\log X_j, \log X_k)
\end{equation}
We compute (Section~\ref{sec:results}): location stress accounts for
27.2\% of variation, efficiency for 22.1\%, peaking for 4.9\%, and
covariance terms together account for 46.7\%. The dominance of the
covariance terms is the empirical signature of compounding harm: low-WUE
facilities are systematically located in low-stress regions and
vice-versa, and the joint signal is much larger than the sum of the
marginal signals. An additive composite would discard this information.

\subsection{Properties}

\begin{proposition}[Monotonicity]
PSWI is strictly non-decreasing in each of its three input factors,
holding the others fixed.
\end{proposition}
\begin{proof}
Immediate from $\partial \textrm{PSWI} / \partial X_k = \prod_{j\neq k} X_j > 0$
for $X_k \in \{\textrm{WUE}, \beta, \textrm{BWS}\}$, since all factors are
non-negative.
\end{proof}

\begin{proposition}[Decomposability]
Cross-facility variation in $\log\textrm{PSWI}$ is exactly the sum of
component variances and covariance terms (Eq.\ above), enabling
attribution of changes to efficiency, peaking, or stress separately.
\end{proposition}

\begin{proposition}[Edge-case correctness]
For facilities where $\textrm{BWS} \to 0$ (sea-water-cooled
non-freshwater-stressed catchments), $\textrm{PSWI} \to 0$, correctly
indicating negligible freshwater-stress impact even when absolute water
volumes are large. For facilities where $\textrm{BWS} > 1$ (extremely
high stress: demand exceeds renewable supply), PSWI amplifies WUE
proportionally, correctly indicating that any consumption in such a
location compounds an unsustainable resource depletion.
\end{proposition}

\subsection{Known Limitations of the Metric}

\begin{tcolorbox}[colback=yellow!8, colframe=yellow!50!black,
                  title=Box 1: Quantified known limitations]
\textbf{Peaking factor estimation.} $\beta_i$ is estimated by climate
zone with documented uncertainty $\pm 30\%$. Sensitivity analysis
(Section~\ref{sec:results}) shows that $\pm 50\%$ per-facility
perturbation produces mean Spearman $\rho = 0.982$ between baseline
and perturbed rankings. \\[2pt]
\textbf{BWS uncertainty.} Aqueduct 4.0 documents catchment-level BWS
uncertainty of approximately $\pm 25\%$ from PCR-GLOBWB 2 model
error~\cite{kuzma2023aqueduct}. Our Monte Carlo counterfactual
propagates this. \\[2pt]
\textbf{Annual BWS averaging.} Aqueduct also publishes monthly BWS
data; switching to monthly values is the highest-priority
extension. \\[2pt]
\textbf{Scope 2 omission.} Direct water (Scope~1) is the focus
because it predominantly draws from public water systems. A Scope~2
robustness check (Section~\ref{sec:results}) shows our findings are
preserved (rank correlation $\rho = 0.946$) when $1.5$~L/kWh of
estimated grid water is added uniformly.
\end{tcolorbox}

% ============================================================================
\section{Methodology}
\label{sec:methodology}
% ============================================================================

\subsection{Sample Frame and Inclusion Criteria}

We include a facility if and only if all of the following hold:
(i)~it operates as either a hyperscale or major colocation data center
(MW-scale IT load); (ii)~its operator publishes either site-specific or
fleet-aggregated WUE for the most recent reporting year (2023 or 2024);
(iii)~its physical location is publicly known to at least the
metropolitan-area level; and (iv)~its catchment is covered by WRI
Aqueduct 4.0. We exclude operators that do not publish any WUE
(e.g.,~xAI was retained because press estimates exist for Memphis), and
operators whose published water data we could not link to specific
facilities (e.g.,~Chinese hyperscalers Alibaba, Tencent, Baidu).

\paragraph{Coverage estimate.}
The resulting sample of 53 facilities represents an estimated 35--45\% of
total U.S. hyperscale data center IT capacity reported in LBNL's 2024
census~\cite{lbnl2024} when weighted by the fleet-share data each
operator discloses. Coverage is best for U.S. and Northern European
hyperscalers; weakest for South / Southeast Asia and entirely absent for
mainland China. We treat the sample as broadly representative of disclosed
hyperscale capacity in the OECD, with explicit caution against
generalization to undisclosed operators.

\subsection{Data Sources}

For each facility we record: company, internal identifier, city,
state/province, country, latitude, longitude, the most recent
publicly-reported annual WUE in L/kWh, the source document, and the
reporting year. Sources include corporate sustainability reports for
2024~\cite{google2025envreport, microsoft2024sustain, meta2025sustain,
amazon2024aws, equinix2024sustain, digitalrealty2024impact,
qts2024sustain}; government records cited in
Han et al.~\cite{han2026}; and LBNL's 2024 Data Center Energy Usage
Report~\cite{lbnl2024} for fleet-share weighting.

\subsection{Water Stress Assignment}

For each facility, we assign a Baseline Water Stress (BWS) value from
WRI Aqueduct 4.0~\cite{kuzma2023aqueduct} at the catchment level. We use
the raw BWS ratio rather than Aqueduct's categorical labels to preserve
quantitative resolution. The complete water-stress assignment table is
released as \texttt{water\_stress.csv}.

\subsection{Peaking Factor Estimation}

Operators rarely disclose peaking factors. Han et al.~\cite{han2026}
document observed $\beta$ of 6--10 for evaporative-cooled facilities in
hot summer climates. We assign $\beta$ by Köppen climate zone:
cool/subarctic 1.8; mild oceanic 2.5; temperate continental 4.0; hot
humid 6.0; hot arid 8.0. Each value is at or below the lower end of Han
et al.'s observed range for similar climates, biasing PSWI estimates
downward. The full sensitivity analysis (Section~\ref{sec:results})
documents that this assignment does not drive our findings.

\subsection{Reproducibility}

All computations are reproduced from the public repository at
\url{https://github.com/hrihaan19/pswi-data-center-audit}. Pipeline
dependencies are limited to \texttt{pandas}, \texttt{numpy},
\texttt{scipy}, and \texttt{matplotlib}. No proprietary or API-gated
data sources are required.

% ============================================================================
\section{Results}
\label{sec:results}
% ============================================================================

\subsection{The Shell Game: WUE vs.\ PSWI Rankings}

\begin{figure}[t]
\centering
\includegraphics[width=\columnwidth]{figure_1_shell_game.png}
\caption{Rank shifts between self-reported WUE and PSWI for the eight
audited companies. Microsoft moves four positions worse under PSWI; xAI
moves five positions better. The rank shift indicates that
location-blind annual reporting introduces systematic distortions in
cross-company sustainability comparisons.}
\label{fig:shellgame}
\end{figure}

Figure~\ref{fig:shellgame} compares the rank order of the eight audited
companies under self-reported WUE and under PSWI. Microsoft moves four
positions worse, driven by its portfolio's concentration in extremely
high-stress catchments (Goodyear and Mesa, AZ; Tuas, SG); xAI moves five
positions better because its single audited facility in Memphis sits on
a moderately stressed basin despite its high estimated WUE.

\paragraph{Bootstrap rank stability (Reviewer point \#5).}
We test whether these company-level shifts are statistically meaningful
under resampling. Within each company, we resample its facilities with
replacement 2{,}000 times, recomputing the company's mean WUE and mean
PSWI at each iteration. Figure~\ref{fig:bootshifts} shows the resulting
distributions.

\begin{figure}[t]
\centering
\includegraphics[width=\columnwidth]{figure_8_bootstrap_shifts.png}
\caption{Bootstrap distribution of company rank shifts (2{,}000
resamples; black bars are 95\% CI). xAI's improvement is statistically
robust (95\% CI excludes zero); Microsoft's, Equinix's, and Digital
Realty's deterioration trends are large in central tendency but their
95\% CIs span zero, reflecting the fundamental statistical limit of
$n = 8$ companies.}
\label{fig:bootshifts}
\end{figure}

The bootstrap analysis is honest about a limit of company-level
rankings: with $n = 8$ companies, only xAI's 5-position improvement
clears a 95\% confidence threshold. Microsoft (mean shift $-2.6$,
median $-3$, 95\% CI $[-5, +1]$), Equinix ($-2.0$, $[-4, 0]$), and
Digital Realty ($-2.2$, $[-4, 0]$) all show large central-tendency
shifts whose CIs nominally include zero. Two interpretations are
possible: (a) at $n = 8$ companies, only the most extreme shifts are
statistically distinguishable; (b) the within-company variance in PSWI
is large because facility portfolios are heterogeneous. We treat the
\emph{facility-level} findings ($n = 53$) as our primary statistical
contribution and the company-level shifts as suggestive directional
evidence.

\subsection{The 4{,}100$\times$ Disparity}

\begin{figure}[t]
\centering
\includegraphics[width=\columnwidth]{figure_2_disparity.png}
\caption{PSWI distribution by company on a logarithmic scale. The best
facility (Hamina, Finland; PSWI $= 0.0036$) and the worst (Phoenix, AZ;
PSWI $= 14.78$) differ by approximately 4{,}100$\times$.}
\label{fig:disparity}
\end{figure}

Figure~\ref{fig:disparity} shows the PSWI distribution by company. The
range spans nearly four orders of magnitude. Google's portfolio is the
most internally diverse, with Henderson, NV (PSWI $= 8.64$) above
Digital Realty's company average and Hamina, FI (PSWI $= 0.0036$) the
lowest in the dataset.

\subsection{Same-Operator Case Study (Reviewer Point \#7)}

To control for reporting methodology, we examine pairs of facilities
operated by the \emph{same company}---thereby holding measurement
practices, internal accounting, and corporate culture
constant---and compare their PSWI scores. Within Google's portfolio:
\begin{itemize}
    \item Lenoir, NC: WUE $= 0.7$~L/kWh, BWS $= 0.18$, $\beta = 4.0$,
        PSWI $= 0.504$
    \item Henderson, NV: WUE $= 0.9$~L/kWh, BWS $= 1.20$, $\beta = 8.0$,
        PSWI $= 8.64$
\end{itemize}
The two facilities report similar self-reported annual WUE
(within 0.2~L/kWh of each other), are operated by the same company under
identical reporting standards, and yet differ in PSWI by a factor of 17.
The driver is location: Lenoir lies in the Catawba River basin
(BWS $= 0.18$); Henderson in the Lower Colorado River basin
(BWS $= 1.20$). This isolated comparison rules out cross-company
methodology differences as an explanation for the dispersion documented
in Figure~\ref{fig:wuepswi}.

\begin{figure}[t]
\centering
\includegraphics[width=\columnwidth]{figure_4_wue_vs_pswi.png}
\caption{Each point is one facility, colored by Köppen climate zone.
The y-axis is logarithmic. Within bands of similar self-reported WUE,
PSWI varies by more than an order of magnitude.}
\label{fig:wuepswi}
\end{figure}

\subsection{Variance Decomposition (Reviewer Point \#8)}

\begin{figure}[t]
\centering
\includegraphics[width=\columnwidth]{figure_9_variance_decomposition.png}
\caption{Variance decomposition of $\log(\textrm{PSWI})$ across 53
facilities. Location stress (BWS) accounts for 27.2\% of variation,
efficiency (WUE) for 22.1\%, peaking ($\beta$) for 4.9\%; covariance
terms contribute 46.7\%, indicating that the components compound rather
than add.}
\label{fig:variance}
\end{figure}

Figure~\ref{fig:variance} shows the variance decomposition. Two
findings:

\textbf{Location stress dominates marginal variation.} BWS (27.2\%) and
WUE (22.1\%) together account for nearly half of cross-facility
variation; the peaking factor (4.9\%) accounts for far less. This is
consistent with our sensitivity analysis: PSWI rankings are robust to
$\beta$ misspecification because $\beta$ contributes relatively little
marginal information.

\textbf{Covariance terms dominate joint variation.} The 46.7\%
covariance share indicates that the three components are not
independent: facilities with high WUE tend to be located in high-stress
regions (and to have higher peaking factors via climate). The
covariance signal is the empirical signature of \emph{compounding}
harm and is the central justification for the multiplicative
functional form.

\subsection{Sensitivity to $\beta$ (Reviewer Point \#1)}

\begin{figure}[t]
\centering
\includegraphics[width=\columnwidth]{figure_7_beta_sensitivity.png}
\caption{(a) Uniform $\beta$ rescaling preserves rankings exactly.
(b) Per-facility independent $\beta$ perturbation of $\pm$50\% across
1{,}000 runs yields mean Spearman $\rho = 0.982$ between baseline and
perturbed rankings, with mean top-10 worst overlap of 7.9 out of 10.}
\label{fig:beta}
\end{figure}

Figure~\ref{fig:beta} reports two tests. Panel (a) tests whether a
constant misspecification of the peaking-factor table preserves
rankings: it does, exactly, since rescaling $\beta$ uniformly is a
monotonic transformation of PSWI. Panel (b) is a stricter test: each
facility's $\beta$ is independently perturbed by a random $\pm 50\%$,
1{,}000 times. The mean Spearman $\rho$ between baseline and perturbed
rankings is 0.982; the 5th percentile is 0.975. The mean top-10 worst
overlap is 7.9 out of 10. Together these establish that PSWI rankings
are robust to plausibly large misspecifications of the peaking factor.

\subsection{Scope 2 Robustness (Reviewer Point \#10)}

We test whether including approximate Scope 2 (off-site
electricity-related) water reverses our findings, using a uniform US
grid average of 1.5~L/kWh from EPRI/NREL benchmarks. Recomputing PSWI
with the augmented WUE yields Spearman $\rho = 0.946$ between original
and Scope-2-inclusive rankings, with 7 of 10 worst-PSWI facilities
preserved. Mean absolute rank change is 3.32 positions. Our findings
are robust to Scope 2 inclusion at the rank level.

\subsection{Geographic Concentration of Impact}

\begin{figure*}[t]
\centering
\includegraphics[width=0.85\textwidth]{figure_3_world_map.png}
\caption{Global distribution of audited facilities, colored by
$\log_{10}(\textrm{PSWI})$. The U.S. Southwest and Singapore/Taiwan
emerge as high-PSWI clusters; Northern Europe and the U.S. Pacific
Northwest are uniformly low.}
\label{fig:map}
\end{figure*}

Figure~\ref{fig:map} maps facilities globally. Of the 10 worst PSWI
facilities, six are in the U.S. Southwest (Phoenix, Henderson, Salt
Lake City, Mesa) and four are in Asia-Pacific (Singapore, Taiwan). Of
the 10 best, seven are in Northern Europe and three in the U.S. Pacific
Northwest. This concentration reflects an unfortunate alignment between
regions attractive to data center operators and regions experiencing
water stress.

\subsection{Disclosure Compliance}

\begin{figure}[t]
\centering
\includegraphics[width=\columnwidth]{figure_6_disclosure.png}
\caption{Disclosure compliance heatmap for major data center operators.
Only Google publishes site-specific WUE; \emph{no} operator publishes
peak or monthly water data, the metric most relevant for public-water-
system stress assessment.}
\label{fig:disclosure}
\end{figure}

Figure~\ref{fig:disclosure} audits the disclosure practices of 12 major
operators against five categories of water data. The picture is bleak:
none publish peak or monthly water data. This is the empirical
foundation of our policy proposal.

% ============================================================================
\section{Counterfactual Analysis and Policy Implications}
\label{sec:policy}
% ============================================================================

\subsection{Counterfactual: PSWI-Aware Reallocation with Full Uncertainty Propagation}

\begin{figure}[t]
\centering
\includegraphics[width=\columnwidth]{figure_10_counterfactual_uncertain.png}
\caption{Counterfactual impact reduction under PSWI-aware workload
reallocation, with documented uncertainty propagated through all three
inputs (WUE $\pm$15\%, BWS $\pm$25\%, $\beta$ $\pm$30\%). At 10\%
fleet reallocation: median reduction 36.7\%, 95\% CI $[26.1\%, 47.0\%]$.}
\label{fig:counterfactual}
\end{figure}

The Monte Carlo counterfactual (Figure~\ref{fig:counterfactual})
propagates documented relative uncertainty through all three PSWI inputs:
WUE $\pm$15\% (corporate reporting precision), BWS $\pm$25\%
(PCR-GLOBWB 2 model uncertainty per Aqueduct 4.0~\cite{kuzma2023aqueduct}),
and $\beta$ $\pm$30\% (Han et al.~\cite{han2026} observed range). At 10\%
fleet workload reallocation, the median impact reduction is 36.7\% with a
95\% CI of $[26.1\%, 47.0\%]$ and a 50\% CI of $[33.3\%, 40.5\%]$. At
20\% reallocation: 60\%, CI $[53\%, 67\%]$; at 30\%: 70\%, CI
$[64\%, 75\%]$. The lower confidence bounds remain substantial, and the
strongly concave shape of the curve is preserved across all 500
simulation runs.

\subsection{Disclosure Framework}

We propose a regulatory framework, modeled on SEC Regulation
S-K~\cite{secregsk}, requiring per-facility, per-quarter PSWI
disclosure. Three core requirements:
(i)~per-facility annual WUE (no aggregation);
(ii)~peak monthly WUE for each facility, enabling computation of a
12-month rolling peak as $\beta$;
(iii)~source-resolved water inputs (potable vs.\ non-potable;
withdrawal vs.\ consumption).

\paragraph{Relation to existing climate disclosure rules.}
The SEC's climate disclosure rules (final rule released March 2024;
partially stayed pending litigation) faced
nearly-identical objections: data quality, third-party assurance, and
small-issuer burden. The SEC addressed these via tiered phase-in
(large accelerated filers first), graduated assurance requirements
(``limited assurance'' before ``reasonable assurance''), and
materiality thresholds. The same architecture can be applied to PSWI
disclosure: hyperscale operators (>50~MW IT load) phase in first;
operators above 5~MW second; below threshold exempt. This pragmatic
phasing preserves regulatory ambition without requiring small-business
exemption negotiation.

\paragraph{Enforcement leverage.}
A novel feature is that enforcement does not require new statutory
authority. U.S. public water systems already have permitting authority
over large water users under the Safe Drinking Water Act and
state-level analogs~\cite{epa2024watercongress}. PSWI disclosure can be
required as a condition of permitting, leveraging existing utility
authority.

% ============================================================================
\section{Limitations and Future Work}
\label{sec:limitations}
% ============================================================================

\paragraph{Annual BWS averaging.}
We use Aqueduct's annual BWS values. Aqueduct also publishes monthly
data; switching to monthly BWS is conceptually consistent with PSWI's
emphasis on summer-peak harm and is the highest-priority extension.

\paragraph{Disclosure-driven coverage.}
Coverage skews toward operators that disclose more data; we have zero
coverage of Chinese hyperscalers and limited coverage of Middle East
and India operators. Adding partial PSWI estimates with explicit
uncertainty flags for these markets is planned.

\paragraph{$\beta$ point-estimation versus measurement.}
Our $\beta$ estimates are climate-zone-conditional, not measurement-based.
Sensitivity analysis (Section~\ref{sec:results}) shows rankings are
robust, but a future revision will use disclosed monthly water data
(currently available only from Google~\cite{google2025envreport}) or
hourly cooling-degree-day modeling.

\paragraph{No causal claim.}
We do not claim PSWI \emph{causes} specific water-system harms.
Empirical validation against observed water-system stress events is a
natural follow-up.

% ============================================================================
\section{Conclusion}
\label{sec:conclusion}
% ============================================================================

The water sustainability of AI data centers is shaped by two factors
that current industry metrics conceal: temporal concentration of
demand and spatial heterogeneity of stress. We have introduced PSWI
to integrate both, applied it to the first cross-company global
dataset of major data centers, and documented robustness-tested
findings: a 4{,}100$\times$ disparity between best and worst facilities;
a same-operator 17$\times$ ratio that controls for reporting
methodology; rank stability under $\pm 50\%$ $\beta$ perturbation;
and a counterfactual showing that 10\% workload reallocation yields a
36.7\% (95\% CI $[26.1, 47.0]$) reduction in fleet impact under
documented uncertainty in all inputs. We have proposed a regulatory
framework engaging directly with the SEC's actual climate disclosure
architecture. All data and code are publicly released.

\section*{Acknowledgments}

The author thanks the authors of Han et al.~(2026), Wu et al.~(2025),
and Li et al.~(2024, 2025) for foundational work that made this
synthesis possible. The World Resources Institute provided the
Aqueduct 4.0 dataset under a Creative Commons license. We thank an
anonymous reviewer whose detailed critique substantially improved this
work.

\bibliographystyle{IEEEtran}
\bibliography{references}

\end{document}
